\begin{description}
\item[(a)] As the number of H-atoms undergoing ionization and recomination % sic
  are balanced at $R_S$, each photon can ionize exactly one hydrogen atom and
  each neutral hydrogen has exactly one proton and one electron,
  \begin{align*}
    n_{\mathrm{recomb}} &= n_{HII} = Q\\
    \text{and } n_e &= n_p = n_H \tag*{(2.0)}\\
    n_{\mathrm{recomb}} &= \alpha n_p n_e V_S\\
    Q &= \alpha n_H^2\,\frac{4\pi}{3}R_S^3 \tag*{(1.0)}\\
    \therefore R_S &= \sqrt[3]{\frac{3Q}{4\pi\alpha n_H^2}} \tag*{(1.0)}\\
      &= \sqrt[3]{\frac{3\times10^{49}}{4\pi\times10^{-19}\times(10^{8})^2}}\\
      &= 1.3\times10^{17}\,\mathrm{m}\\
    \therefore R_S &\approx 4\,\mathrm{pc} \tag*{(1.0)}
  \end{align*}
\item[(b)] Per one unit length distance, a typical photon will encounter
  $\sigma n_H$ H-atoms. Thus, the mean free path will be,
  \begin{align*}
    l_{\nu_0} &= \frac{1}{\sigma n_H} = \frac{1}{\sigma 10^{-21}\times10^{8}}
    % sic: the source prints "\sigma 10^{-21}" in the denominator
    \\
    l_{\nu_0} &= 10^{13}\,\mathrm{m} \tag*{(2.0)}\\
    \therefore l_{\nu_0} &\approx 10^{-4}R_S \ll R_S
  \end{align*}
  Thus, the boundary layer of the sphere is very thin as compared to its total
  size. Hence, ``YES'' this ionized nebula is very sharp-edged. \hfill(1.0)
\item[(c)] For Stromgren sphere to form, all H-atoms ($N$) inside $R_S$ need
  to be ionized. Thus, time $t_S$ required will be
  \begin{align*}
    N &= V_S n_H = \frac{4\pi}{3}R_S^3 n_H \tag*{(1.0)}\\
    t_S &= \frac{N}{Q} = \frac{4\pi R_S^3 n_H}{3Q} \tag*{(1.0)}\\
        &= \frac{1}{\alpha n_H} = \frac{1}{10^{-19}\times10^{8}}\\
        &= 10^{11}\,\mathrm{s}\\
        &\approx 3000\,\mathrm{yr} \tag*{(2.0)}
  \end{align*}
  For typical nebular densities, the main-sequence lifetime of LBV stars is
  much longer than this ionization time, and hence our assumption of a
  stationary system is justified.
\item[(d)] The mean electron thermal energy in a plasma of temperature
  $T_e = 10^4$\,K is
  \[
    E_e \approx k_B T_e \approx 1.4\times10^{-19}\,\mathrm{J}
        \approx 0.9\,\mathrm{eV}
  \]
  \hfill(2.0)

  The energy of a photon is
  \[
    E_\gamma = h\nu \approx 6.6\times10^{-24}\,\mathrm{J}
             \approx 4\times10^{-5}\,\mathrm{eV}
  \]
  \hfill(1.0)

  As the energy of the radio photon is much much smaller than that of the
  electron, the answer is ``Yes''. \hfill(1.0)
\item[(e)] As the plasma is nearly transparent, using the Rayleigh--Jeans law,
  \begin{align*}
    B_\nu &= \frac{2k_B T_e \nu^2}{c^2} \tag*{(1.0)}\\
    S_\nu &\propto B_\nu \tau_\nu = \frac{2k_B T_e \nu^2}{c^2}\,\tau_\nu\\
    S_\nu &\propto \nu^2\cdot\nu^{-2.1}\\
    \therefore S_\nu &\propto \nu^{-0.1} \tag*{(3.0)}
  \end{align*}
\end{description}
